% === File 0105.md (Printed Page 103: ಷಟ್ಪದಿಗಳ ಲಕ್ಷಣಗಳು, full-width upright table) ===
\par\vspace{6pt plus 4pt}\noindent\begin{minipage}{\textwidth}% keep title and table together
\begin{center}{\kbvchapfont\fontsize{14pt}{18pt}\selectfont ಷಟ್ಪದಿಗಳ ಲಕ್ಷಣಗಳು}\end{center}
\vspace{4pt plus 3pt minus 1pt}
\noindent{\fontsize{8.8pt}{11pt}\selectfont\renewcommand{\arraystretch}{1.05}\setlength{\tabcolsep}{3pt}\arrayrulecolor{kbvline}%
\setlength{\kbvcolw}{\dimexpr(\textwidth-17mm-21mm-6\tabcolsep)/2\relax}%
\begin{tabular}{@{} >{\raggedright\arraybackslash}p{17mm} >{\raggedright\arraybackslash}p{\kbvcolw} >{\raggedright\arraybackslash}p{\kbvcolw} >{\raggedright\arraybackslash}p{21mm} @{}}
      \specialrule{0.5pt}{0pt}{3pt}
      {\kbvsecfont ಷಟ್ಪದಿ} & {\kbvsecfont ೧, ೨, ೪, ೫ ನೆಯ ಪಾದಗಳು} & {\kbvsecfont ೩ ಮತ್ತು ೬ ನೆಯ ಪಾದಗಳು} & {\kbvsecfont ವಿಶೇಷ ವಿಚಾರ} \\
      \specialrule{0.5pt}{3pt}{3pt}
      {\kbvsecfont ೧.}\,ಶರ ಷಟ್ಪದಿ & ಪ್ರತಿ ಪಾದದಲ್ಲಿ ನಾಲ್ಕು ಮಾತ್ರೆಯ ಎರಡು ಗಣಗಳಿರುವುವು. & ಪ್ರತಿ ಪಾದದಲ್ಲಿ 4 ಮಾತ್ರೆಯ 3 ಗಣಗಳು ಮತ್ತು ಮೇಲೊಂದು ಗುರು ಇರುವುದು. & “\laghu\,\guru\,\laghu” ಈ ರೀತಿಯ ಗಣ ಬರುವುದಿಲ್ಲ. \\
      \specialrule{0.3pt}{2.5pt}{2.5pt}
      {\kbvsecfont ೨.}\,ಕುಸುಮ ಷಟ್ಪದಿ & ಪ್ರತಿ ಪಾದದಲ್ಲಿ ಐದು ಮಾತ್ರೆಯ ಎರಡು ಗಣಗಳಿರುವುವು. & ಪ್ರತಿ ಪಾದದಲ್ಲಿ 5 ಮಾತ್ರೆಯ 3 ಗಣಗಳು ಮತ್ತು ಮೇಲೊಂದು ಗುರು ಇರುವುದು. & “\laghu\,\guru\,\laghu\laghu” ಮತ್ತು “\laghu\,\guru\,\guru” ರೀತಿಯ ಗಣವಿರುವುದಿಲ್ಲ. \\
      \specialrule{0.3pt}{2.5pt}{2.5pt}
      {\kbvsecfont ೩.}\,ಭೋಗ ಷಟ್ಪದಿ & ಪ್ರತಿ ಪಾದದಲ್ಲಿ ಮೂರು ಮಾತ್ರೆಯ ನಾಲ್ಕು ಗಣಗಳಿರುವುವು. & ಪ್ರತಿ ಪಾದದಲ್ಲಿ 3 ಮಾತ್ರೆಯ 6 ಗಣಗಳಿರುವುವು ಮತ್ತು ಮೇಲೊಂದು ಗುರುವಿರುವುದು. & “\laghu\,\guru” ಈ ರೀತಿಯ ಗಣಗಳಿರುವುದಿಲ್ಲ. \\
      \specialrule{0.3pt}{2.5pt}{2.5pt}
      {\kbvsecfont ೪.}\,ಭಾಮಿನೀ ಷಟ್ಪದಿ & ಪ್ರತಿ ಪಾದದಲ್ಲಿ 3 ಮತ್ತು 4 ಮಾತ್ರೆಯ 4 ಗಣಗಳು ಪರ್ಯಾಯವಾಗಿ ಬರುವುವು. & ಪ್ರತಿ ಪಾದದಲ್ಲಿ 3 ಮತ್ತು 4 ಮಾತ್ರೆಯ 6 ಗಣಗಳು ಪರ್ಯಾಯವಾಗಿ ಬರುವುವು ಮತ್ತು ಮೇಲೊಂದು ಗುರುವಿರುವುದು. & “\laghu\,\guru\,\laghu” ಈ ರೀತಿಯ ಗಣವಿರುವುದಿಲ್ಲ. \\
      \specialrule{0.3pt}{2.5pt}{2.5pt}
      {\kbvsecfont ೫.}\,ಪರಿವರ್ಧಿನೀ ಷಟ್ಪದಿ & ಪ್ರತಿ ಪಾದದಲ್ಲಿ 4 ಮಾತ್ರೆಯ 4 ಗಣಗಳಿರುವುವು. & ಪ್ರತಿ ಪಾದದಲ್ಲಿ 4 ಮಾತ್ರೆಯ 6 ಗಣಗಳಿರುವುವು ಮತ್ತು ಮೇಲೊಂದು ಗುರುವಿರುವುದು. & “\laghu\,\guru\,\laghu” ಈ ರೀತಿಯ ಗಣವಿರುವುದಿಲ್ಲ. \\
      \specialrule{0.3pt}{2.5pt}{2.5pt}
      {\kbvsecfont ೬.}\,ವಾರ್ಧಕ ಷಟ್ಪದಿ & ಪ್ರತಿ ಪಾದದಲ್ಲಿ 5 ಮಾತ್ರೆಯ 4 ಗಣಗಳಿರುವುವು. & ಪ್ರತಿ ಪಾದದಲ್ಲಿ 5 ಮಾತ್ರೆಯ 6 ಗಣಗಳು ಮತ್ತು ಮೇಲೊಂದು ಗುರು ಇರುವುದು. & {\color{kbvgrey}\ldots\ldots} \\
      \specialrule{0.5pt}{3pt}{0pt}
\end{tabular}\arrayrulecolor{black}}
\end{minipage}\par\vspace{6pt plus 4pt}
